\documentclass[11pt]{article}

\usepackage[margin=1in]{geometry}
\usepackage{amsmath,amssymb,amsthm,mathtools}
\usepackage{lmodern}
\usepackage[T1]{fontenc}

\title{Two-Type Example: Capacity Sorting in the $\beta=0$ Regime}
\date{}

\begin{document}
\maketitle

Consider a two-type version of the within-category problem, with a low type $\theta_L$, a high type $\theta_H>\theta_L$, and population shares $p\in(0,1)$ and $1-p$. Let the social weights be $\lambda_L$ and $\lambda_H$, and define the average weight
\[
\tilde{\lambda}=p\lambda_L+(1-p)\lambda_H.
\]
We focus on the regime $\beta=0$, so that the effective weights coincide with $\Gamma_L$ and $\Gamma_H$.

\section*{1. Two-type objects}

The two-type analogue of the revenue-productivity term is
\[
J_L=\theta_L-\frac{1-p}{p}\left(\theta_H-\theta_L\right),
\qquad
J_H=\theta_H.
\]
The corresponding effective welfare weights are
\[
\Gamma_L=\tilde{\lambda}J_L+\Lambda_L,
\qquad
\Gamma_H=\tilde{\lambda}J_H,
\]
with
\[
\Lambda_L=\frac{1-p}{p}\left(\theta_H-\theta_L\right)\lambda_H,
\qquad
\Lambda_H=0.
\]
Equivalently,
\[
\Gamma_H=\tilde{\lambda}\,\theta_H,
\]

and
\[
\Gamma_L
=
\tilde{\lambda}\left[\theta_L-\frac{1-p}{p}\left(\theta_H-\theta_L\right)\right]
+
\frac{1-p}{p}\left(\theta_H-\theta_L\right)\lambda_H.
\]
A useful alternative expression is
\[
\Gamma_L
=
\tilde{\lambda}\,\theta_L
+
\frac{\theta_H-\theta_L}{p}\left(\lambda_H-\tilde{\lambda}\right).
\]

\section*{2. Allocation and capacity expansion when $\beta=0$}

Let $q_L$ and $q_H$ denote the quantities assigned to the low and high type in a binding state. The first-order conditions are
\[
\Gamma_L\,u(q_L)=\varepsilon,
\qquad
\Gamma_H\,u(q_H)=\varepsilon,
\]
where $u=U_q$ is marginal utility and $\varepsilon$ is the scarcity rent. The capacity constraint is
\[
pq_L+(1-p)q_H=k.
\]

Differentiating with respect to $k$ gives
\[
\Gamma_Lu_q(q_L)\frac{dq_L}{dk}=\frac{d\varepsilon}{dk},
\qquad
\Gamma_Hu_q(q_H)\frac{dq_H}{dk}=\frac{d\varepsilon}{dk},
\qquad
p\frac{dq_L}{dk}+(1-p)\frac{dq_H}{dk}=1.
\]
Therefore,
\[
\frac{dq_L}{dk}
=
\frac{\Gamma_Hu_q(q_H)}{(1-p)\Gamma_Lu_q(q_L)+p\Gamma_Hu_q(q_H)},
\qquad
\frac{dq_H}{dk}
=
\frac{\Gamma_Lu_q(q_L)}{(1-p)\Gamma_Lu_q(q_L)+p\Gamma_Hu_q(q_H)}.
\]

Define
\[
R_{dq}:=\frac{dq_H/dk}{dq_L/dk}.
\]
Then
\[
R_{dq}=\frac{w_H}{w_L},
\qquad
w_L:=-\frac{1}{\Gamma_Lu_q(q_L)},
\qquad
w_H:=-\frac{1}{\Gamma_Hu_q(q_H)}.
\]
Using the first-order conditions,
\[
\frac{\Gamma_L}{\Gamma_H}=\frac{u(q_H)}{u(q_L)},
\]
so
\[
R_{dq}=\frac{u(q_H)}{u(q_L)}\frac{u_q(q_L)}{u_q(q_H)}.
\]
If we define the curvature object
\[
R_q(q):=-\frac{u(q)}{u_q(q)},
\]
then
\[
R_{dq}=\frac{R_q(q_H)}{R_q(q_L)}.
\]
Thus, outside the CARA case, a marginal capacity increase need not be split one-for-one across types; the split is governed by the ratio of the curvature terms.

\section*{3. The capacity-sorting object}

The discrete analogue of the revenue-productivity term is
\[
EUJ=pJ_LU(q_L)+(1-p)J_HU(q_H).
\]
Replacing $J_L$ and $J_H$ yields
\[
EUJ
=
\left[p\theta_L-(1-p)\left(\theta_H-\theta_L\right)\right]U(q_L)
+
(1-p)\theta_HU(q_H).
\]
Differentiating with respect to $k$ gives
\[
\frac{\partial EUJ}{\partial k}
=
\left[p\theta_L-(1-p)\left(\theta_H-\theta_L\right)\right]u(q_L)\frac{dq_L}{dk}
+
(1-p)\theta_Hu(q_H)\frac{dq_H}{dk}.
\]
Factorizing by $dq_L/dk$ and using $R_{dq}=\left(dq_H/dk\right)/\left(dq_L/dk\right)$,
\[
\frac{\partial EUJ}{\partial k}
=
\frac{dq_L}{dk}
\left\{
\theta_Lu(q_L)
-
(1-p)\theta_H\left[u(q_L)-R_{dq}u(q_H)\right]
\right\}.
\]
Replacing marginal utility by the first-order conditions and factorizing by $\varepsilon$ gives
\[
\frac{\partial EUJ}{\partial k}
=
\frac{dq_L}{dk}\,\varepsilon
\left[
\frac{\theta_L}{\Gamma_L}
-
(1-p)\theta_H
\left(
\frac{1}{\Gamma_L}
-
\frac{R_{dq}}{\Gamma_H}
\right)
\right].
\]
Finally, using
\[
R_{dq}=\frac{R_q(q_H)}{R_q(q_L)},
\]
one obtains
\[
\frac{\partial EUJ}{\partial k}
=
\frac{dq_L}{dk}\,\varepsilon
\left[
\frac{\theta_L}{\Gamma_L}
-
(1-p)\theta_H
\left(
\frac{1}{\Gamma_L}
-
\frac{1}{\Gamma_H}\frac{R_q(q_H)}{R_q(q_L)}
\right)
\right].
\]
This is the key equation for the two-type example. It isolates three forces:
\begin{enumerate}
    \item the direct low-type contribution $\theta_L/\Gamma_L$;
    \item the standard high-type rent-pressure term $(1-p)\theta_H$;
    \item the peak-load allocation margin, through the curvature ratio $R_q(q_H)/R_q(q_L)$.
\end{enumerate}

\section*{4. CARA benchmark}

Under CARA, $R_q(q)$ is constant, so
\[
\frac{R_q(q_H)}{R_q(q_L)}=1.
\]
Hence
\[
\frac{dq_L}{dk}=1,
\qquad
\frac{dq_H}{dk}=1,
\]
and the previous expression simplifies to
\[
\frac{\partial EUJ}{\partial k}
=
\varepsilon
\left[
\frac{\theta_L}{\Gamma_L}
-
(1-p)\theta_H\left(\frac{1}{\Gamma_L}-\frac{1}{\Gamma_H}\right)
\right].
\]
Setting this equal to zero yields
\[
\theta_L\Gamma_H=(1-p)\theta_H\left(\Gamma_H-\Gamma_L\right).
\]
Using
\[
\Gamma_H=\tilde{\lambda}\,\theta_H,
\qquad
\Gamma_H-\Gamma_L=\lambda_L\left(\theta_H-\theta_L\right),
\]
the cutoff is
\[
\lambda_L^{\ast}
=
\frac{\theta_L\tilde{\lambda}}{(1-p)\left(\theta_H-\theta_L\right)}.
\]
Therefore,
\[
\frac{\partial EUJ}{\partial k}<0
\qquad\Longleftrightarrow\qquad
\lambda_L>\lambda_L^{\ast}.
\]
If one imposes
\[
\tilde{\lambda}=1,
\qquad
\theta_H=\theta_L+1,
\]
then
\[
\lambda_L^{\ast}=\frac{\theta_L}{1-p}.
\]

\section*{5. Economic interpretation}

The two-type example makes the capacity-sorting tension transparent. The negative force does not come from redistribution itself, but from the fact that the low type may have weak or even negative revenue productivity through $J_L$. Redistribution keeps the low type in the allocation through a large effective weight $\Gamma_L$. A capacity expansion can therefore lower $EUJ$ when the planner puts enough weight on a type whose utility is difficult to finance under incentive compatibility.

The CARA cutoff makes this tension especially clear:
\[
\lambda_L^{\ast}
=
\frac{\theta_L\tilde{\lambda}}{(1-p)\left(\theta_H-\theta_L\right)}.
\]
A larger $\theta_L$ raises the cutoff, making it harder to obtain $\partial EUJ/\partial k<0$. Intuitively, when the low type becomes less revenue-destroying, a stronger redistributive tilt is needed for a capacity expansion to reduce the revenue-productivity term.

\end{document}
